% !TeX root = ../概率论与数理统计笔记.tex
% 由主文档通过 \input 载入；请编译主文档。
\setcounter{chapter}{17} % 保留原书章号。
\chapter{贝叶斯统计初步（选读）}
\label{chap:ps-bayesian}

本章在统计模型中为未知参数引入概率分布，研究如何由先验与数据得到后验推断。
先修内容为贝叶斯公式、条件期望与似然；学习目标是理解推断所依赖的先验、模型和损失函数。

\section{先验、似然与后验}
\label{sec:ps-bayesian-model}
定义\emph{先验分布}（prior distribution）、\emph{后验分布}（posterior distribution）与边缘似然（marginal likelihood）。
在连续参数情形，用先验密度与似然的乘积构造后验，并检查归一化常数有限且非零。
说明这里关于参数的概率陈述属于所建立的贝叶斯模型，与频率学派固定参数的重复抽样解释不同。

\section{共轭模型与后验更新}
\label{sec:ps-conjugate-priors}
介绍\emph{共轭先验}（conjugate prior），推导贝塔--二项、伽马--泊松与方差已知的正态均值模型。
逐项注明超参数、数据统计量和更新后的参数，沿用前文伽马分布的率参数约定。
以顺序更新说明在模型条件一致时批量和逐批处理数据的联系。

\section{贝叶斯估计与可信区间}
\label{sec:ps-bayesian-decisions}
在后验期望损失最小的准则下介绍\emph{贝叶斯估计}（Bayes estimator）；平方损失和绝对损失分别联系后验均值与中位数。
介绍\emph{可信区间}（credible interval），并与第\ref{chap:ps-intervals}章的置信区间比较其概率解释。
最大后验估计单独说明为密度最大化方法，并讨论它依赖参数化的特点。

\section{后验预测与先验敏感性}
\label{sec:ps-posterior-predictive}
通过对参数的后验分布积分构造\emph{后验预测分布}（posterior predictive distribution），同时反映参数不确定性与新观测波动。
比较不同合理先验下的推断，说明不适当先验即使产生形式表达式也需检查后验是否为概率分布。
模型比较与贝叶斯因子仅作进一步学习入口，标明对先验和归一化常数的依赖。

\section{例题、习题与学习检查}
规划伯努利成功概率的顺序更新、泊松强度预测、正态均值估计的收缩现象。
习题要求分别报告先验、似然、后验、损失准则与预测对象。
学习检查避免把不同统计框架中的区间解释混用，并比较数据量与先验影响之间的关系。






